\chapter{Languages that make interpretation inspectable}
\label{ch:languages}

\section{Choose the question before choosing the language}
\label{sec:languages-question}
The gift example needs a condition over classified facts. A consent example
needs a history of grants and withdrawals. A payment agreement may need parties,
assets and deadlines. A dispute about the meaning of an exception needs
alternative readings and their supporting reasons. These are related questions,
but a representation designed for one does not automatically solve the others.
The purpose of reviewing formal legal languages is to identify which structures
help make each question precise.

A formal legal language is a language in which statements about a legal or
normative domain have defined syntax and computational meaning. The word
``legal'' identifies the intended subject matter and useful abstractions. It
does not mean the language has legal authority, that a program written in it
is enforceable, or that translation into it establishes a correct reading.
Authority still comes from the relevant source and institutional decision.

For our product, a useful language must support at least three forms of
inspection. A reviewer needs to see how a condition corresponds to source
text. An engineer needs to know exactly what values and events the program
expects. A verification tool needs unambiguous semantics. A language that is
expressive enough to represent any program may make verification harder; a
language that permits only simple decision tables may fail to represent
important temporal or exception structure.

This chapter develops the relevant mechanisms through examples. It does not
choose a universal replacement for the bank's Java system. The recommended
architecture keeps a bounded executable representation, a richer interpretation
record and explicit adapters for specialized analyses. The relationship between
those representations must itself be checked. Adding more languages without
defining that relationship would create new translation boundaries rather than
redundancy.

\section{Rules do not all behave like ordinary implications}
\label{sec:languages-norms}
An ordinary material implication says that if a proposition is true, another
proposition is true. A legal norm often says that, under specified circumstances,
someone ought to do something. The existence of the duty does not establish
that the action occurred. A rule requiring a document to be supplied cannot
be used as evidence that the document was supplied. The distinction between
facts and normative outputs is central to input/output logics
\citep[secs.~2--4]{inputoutput2000}.

For a simple illustration, suppose an adopted specification says that a
distributor must provide a named warning when a condition holds. The condition
can create a duty record. Performance requires separate evidence of the
warning being delivered to the correct recipient in the required context.
If the program inserts the duty's content into its factual database as though
it had happened, every obligation can appear automatically satisfied.

Prohibitions and permissions also need care. Failure to derive a prohibition
may mean the act is permitted under one explicitly defined logic; in a bank
workflow it may instead mean the relevant sources or facts are missing. The
literature distinguishes strong and weak permissions under particular
defeasible-logical definitions \citep[secs.~2--4]{permissions2012}. Those
definitions are not a general license to treat an empty query result as
authorization. The selected RuleIR gift result remains a subcondition, with
permission supplied only by a separate approved decision policy.

Many legal rules are defeasible: they normally support a conclusion, but an
exception or stronger opposing rule can override that support. A defeasible
rule is not the same as an ordinary implication with an informal comment
attached. Its evaluation must account for applicable opposing rules and
priorities. Defeasible-logic systems formalize different ways to do this,
including rules that support conclusions and defeaters that block opposing
support \citep[sec.~2]{defeasible2000}.

The engineering implication is simple but demanding. ``Has exceptions'' is
not a complete semantics. The specification must say what happens when no
exception is known, one is established, two compete, or an exception remains
unknown. It must also say whether a priority is supplied by legal authority,
an explicit interpretation decision or a programming convention. File order
is not an adequate hidden source of legal priority.

The early logic-program treatment of the British Nationality Act demonstrated
how a selected body of legislation could be decomposed into rules with
inspectable reasoning \citep{sergot1986}. That work is valuable here as an
example of disciplined formalization, including attention to time and negative
information. It should not be read as a claim that ordinary legal prose can
always be converted automatically into a complete logic program.

\section{Controlled English: readable structure with restricted freedom}
\label{sec:languages-logicalenglish}
Ordinary English permits pronouns, elliptical references, implied subjects and
several possible attachments of a qualification. Those features make writing
compact, but they leave a translator with choices. Controlled English restricts
some of that freedom so a sentence can be mapped into logic by explicit rules.
Logical English develops this approach for law and education
\citep[language specification and examples]{logicalenglish2023}.

Consider the sentence ``If it is a gift and is linked to the product or its
type, it is prohibited unless discounted.'' The sentence is short, but ``it''
could refer to the offer or component, ``its'' needs a binding, and ``discounted''
does not identify the exception. A controlled rendering should repeat the
subject where necessary and name each definition. Repetition in the language
is useful if it removes a legally consequential ambiguity.

One possible authoring form for our own restricted profile is:
\begin{lstlisting}[caption={Illustrative LegalMath authoring notation. This is a proposed template, not claimed to be valid Logical English or Catala source.}]
subject: component C in offer O
scope:
  actor(O) satisfies distributor_profile
  activity(O) satisfies authorised_fund_promotion_profile
condition:
  C satisfies gift_definition
  AND (
    C has specific_product_link
    OR C has product_type_link
  )
  AND NOT (C satisfies fee_discount_definition)
consequence:
  selected_gift_restriction_triggered
\end{lstlisting}
The form makes four decisions inspectable: the subject, the scope, the
condition and the named consequence. Every definition identifier refers to
a versioned record. A reviewer can ask why a definition applies without
reverse-engineering the syntax tree. An engineer can translate the connectives
deterministically once the template is fixed.

Variable binding is a practical source of error. If one occurrence of $C$
refers to the voucher and another to the whole offer, the sentence has changed
its subject midway. The authoring language should enforce a declared variable
and type environment. A reference to an undeclared subject should be rejected.
An association such as ``component C belongs to offer O'' should be explicit
when the rule uses facts about both.

Negation requires similar precision. ``There is no evidence that C is a gift''
is not the same assertion as ``evidence establishes that C is not a gift.''
Logical-programming systems may distinguish explicit negation from failure to
prove a proposition. Our evidence model records known false, unknown and
conflict separately. A controlled sentence that asks about missing evidence
should use that information state explicitly; it should not rely on a model
to guess which sense of negation the author intended.

Controlled English improves the boundary between reviewed meaning and
implementation. It does not remove the earlier task of choosing the controlled
sentence from the original circular. A reviewer can still choose the wrong
definition or omit a phrase. That is why the product retains source anchors,
alternative templates and contrasting scenarios alongside the readable rule.
The language makes a chosen interpretation precise enough to challenge.

\section{Catala: defaults, exceptions and source-linked authoring}
\label{sec:languages-catala}
Catala was designed around a recurring pattern in legislation: a general rule
defines a value, and explicit exceptions may replace it. Its core calculus
gives that structure a formal meaning; its authoring approach keeps legal text
near the corresponding code \citep[secs.~2--4]{catala2021}. For a reviewer,
this can be more intelligible than a long sequence of ordinary conditional
statements whose priority is implicit in their order.

A hypothetical policy might normally require an explanation and waive a
particular step under a reviewed condition, while restoring it if the client
requests assistance. The important structure is the relationship between
the general rule and its exceptions. A nested exception can express a priority
that differs from two sibling exceptions. Moving a condition one level in
that structure can change the result even when every word remains present.

For a small teaching example, suppose a base value is false and one exception
returns true when $g$ holds. If $g$ is established, the exception supplies the
value. If $g$ is refuted, the base applies. If the specification supplies two
sibling exceptions and both are active, the system needs a conflict rule.
Catala's core treatment distinguishes an absent applicable value from a
conflict between exceptions; more than one active exception can conflict even
when the values agree. Agreement of the values does not prove that their
legal priority has been resolved.

LegalMath borrows explicit exception structure but adds its own treatment of
incomplete evidence. In the present RuleIR profile, every sibling guard is
evaluated. Two true guards conflict. One true guard with an unknown competitor
remains unknown because the competitor may become active. One true guard with
all competitors false selects its value. When all guards are false, the base
is selected. This is a separately specified design; it is not a claim that
RuleIR implements the whole Catala calculus.

The distinction matters for ensemble comparisons. A Catala model and a RuleIR
model can be useful independent representations, but they do not become
equivalent merely because both have an ``exception'' construct. The adapter
must define how missing evidence, no applicable value and competing exceptions
are represented. A comparison outside that common fragment should return
unsupported or expose the difference, not choose a convenient interpretation.

Catala's compiler work also illustrates how formal claims should be scoped.
The paper establishes properties for a core translation; the production
compiler and surrounding tooling are separate implementation objects. The
retained repository inspection also records an admitted supporting lemma in
the formalization notes. This project has not replayed those proofs. We borrow
the explicit language structure and preservation discipline, while requiring
independent evidence for any concrete adapter used in the bank product.

\citet[sec.~4.2]{huttner2022} discuss legal/programmer collaboration as a way
to keep source meaning and executable rules connected. The useful lesson is
organizational as well as technical. A legal reviewer can challenge an
exception's interpretation while an engineer checks that its computational
consequence matches that interpretation. Neither role should be expected to
certify the other's entire domain from a single code review.

\section{Metadata is part of the meaning: LegalRuleML}
\label{sec:languages-legalruleml}
A formula alone does not say which provision it interprets, which jurisdiction
it concerns, when it applies, or whether a competing formula is also
defensible. LegalRuleML addresses representation of legal rules together with
such contextual and provenance information. The work on legal interpretations
explicitly represents alternative readings associated with source material
\citep[secs.~2--4]{legalinterpretations2014}; the later OASIS specification
provides a broader standardized vocabulary \citep{legalruleml2021}.

For the fee-discount issue, two candidate rules can point to the same source
span while naming different interpretation records. One record might restrict
the exception to an initial reduction; another might include a documented
refund. Their source identity is shared, but their assumptions and consequences
differ. The system should not assign a single interpretation identifier to
both merely because the source paragraph is the same.

The representation must also preserve the authority of a selection. A legal
reviewer may adopt one interpretation for a defined product and date range.
That selection is an event with an author, reason and scope. It does not erase
the alternative or establish that the alternative was syntactically invalid.
A future source clarification may require reconsidering the choice. The
provenance record is what lets the system determine which decisions depend on it.

LegalRuleML is a representation standard, not a turnkey engine for resolving
every legal dispute. A serialized permission or priority still requires a
specified evaluation semantics. Different reasoning engines may support
different fragments or make different assumptions. For the prototype we should
borrow the discipline of explicit context and alternatives without requiring
the entire standard to become the first runtime dependency.

This leads to a useful separation in the proposed storage model. A source span
identifies words. An interpretation claim identifies an adopted or proposed
meaning. A formal rule identifies a computation. A review decision identifies
who accepted which claim under which context. Keeping those identities
separate prevents a hash of the source text from being used as if it also
identified the meaning chosen from that text.

\section{Stipula: a contract as parties, resources and transitions}
\label{sec:languages-stipula}
Stipula becomes relevant when a legal arrangement changes through actions over
time. Its original paper describes agreements, parties, fields, assets,
functions, states and scheduled events \citep[secs.~3--6]{stipula2021}. A
function can be available only to a specified party in a specified state, with
a precondition on the current data. Executing it changes fields or resources
and moves the contract to another state. This is a much richer object than
a single decision-table row.

To understand the mechanism, consider a synthetic service agreement. It begins
in a state awaiting agreement on terms. After the designated parties agree,
it enters an active state. An authorized party can request a permitted service
only while the agreement is active. A withdrawal changes the state so that
the same request is no longer available under that agreement. A timed event
may move the agreement into an expired state. This example illustrates a state
model; it is not a claim about the legal effect of any particular bank consent.

The state matters because the same action name can have different consequences
at different points in the history. A function called \texttt{requestService}
cannot be understood by its input parameters alone if its availability depends
on whether consent is active. A current Boolean copied from a stale database
may also be inadequate. The state should be derived from an attributable
history or a reviewed snapshot plus a complete subsequent event stream.

Stipula's agreement construct explicitly identifies who participates and who
must agree on the initial terms. This is useful as a modeling discipline:
the identity of the party agreeing is not an untyped annotation on a field.
For LegalMath, the analogous review question is who has authority to approve a
meaning, data mapping or release. We borrow the importance of explicit parties
and transitions, not a claim that a regulatory circular is itself a private
contract between the bank and regulator.

Assets in Stipula are distinguished from ordinary fields. An ordinary numeric
field can be copied as information; an asset transfer must account for the
resource leaving one location and arriving at another. In the original
semantics, asset movement prevents creating a negative quantity and preserves
the relevant resource through transfer. This supports questions about duplicate
spending, locked resources and the availability of contractual actions. It is
not a model of market liquidity or a substitute for the bank's general ledger.

For a synthetic escrow, suppose a contract holds a resource amount $b$ and
transfers $q$ to a recipient. A local arithmetic description is
\begin{equation}
 b'=b-q,\qquad r'=r+q,\qquad 0\leq q\leq b,
 \label{eq:asset-transfer}
\end{equation}
where $r$ is the recipient's modeled holding. Adding the first two equalities
gives $b'+r'=b+r$. This is a conservation property of the modeled transfer.
It says nothing about whether the underlying asset exists off-system or whether
the transfer is legally authorized. Those are separate premises. The example
shows how a formal language can support a useful exact property without
claiming to settle the whole legal arrangement.

Scheduled events add another dimension. A Stipula event can schedule a
continuation at a time and execute it only if the contract remains in a
specified state. An action taken before the event can change whether that
continuation is applicable. The ordering of an event and a function call at
the same time is consequently significant. The inspected original semantics
gives scheduled events precedence in the relevant equal-time situation. A bank
adapter must not assume that this is the legally correct ordering for its own
deadlines. The source, clock precision and event-order authority need review.

This example also distinguishes observation from permission. A formal contract
may describe which function calls its own execution platform permits. A bank
audit system must still record actions that occurred outside a permitted
workflow. Deleting an observed transaction because the policy would have
forbidden it would erase evidence of a possible violation. A normative monitor
therefore needs a clear separation between advice about an action and recording
its actual occurrence.

\section{A small state contract that a reader can execute by hand}

A complete miniature model helps distinguish the legal-language idea from its
implementation in a particular tool. Consider a synthetic escrow agreement in
which a payer deposits a fixed amount, a recipient may receive it after an
accepted delivery, and the payer can recover it after an expiry if delivery has
not been accepted. The example is chosen to teach states, parties, resources and
time. It is not a model of an actual bank product or a proposed interpretation of
the gift circular.

Let the contract state be one of \texttt{UNFUNDED}, \texttt{HELD},
\texttt{PAID} and \texttt{REFUNDED}. Let $b$ be the amount held by the contract,
$p$ the payer's modeled external holding and $r$ the recipient's holding. All
amounts are nonnegative integer units. Let $d$ be the agreed deposit, $\tau$ the
expiry time and $t$ the current event time. The parties and the meaning of
accepted delivery are agreed inputs to this model.

The state transition table is the executable meaning of the miniature contract.
It says which caller can perform which action, under what guard, and how the
state and resources change. It is deliberately written in a neutral notation;
it is not presented as tested Stipula source code.

\begin{longtable}{@{}P{24mm}P{33mm}P{42mm}P{48mm}@{}}
\caption{Synthetic escrow transitions, inspired by state-contract languages.}\\
\toprule
From & Action and caller & Guard & Effect and next state\\\midrule\endfirsthead
\toprule From & Action and caller & Guard & Effect and next state\\\midrule\endhead
\texttt{UNFUNDED} & Deposit by payer & $p\geq d$, $d>0$, $t<\tau$ & $p'=p-d$, $b'=d$, $r'=r$; enter \texttt{HELD}.\\
\texttt{HELD} & Release by authorized acceptance process & Accepted delivery and $t<\tau$ & $r'=r+b$, $b'=0$, $p'=p$; enter \texttt{PAID}.\\
\texttt{HELD} & Expiry event & $t\geq\tau$ & $p'=p+b$, $b'=0$, $r'=r$; enter \texttt{REFUNDED}.\\
\texttt{PAID} or \texttt{REFUNDED} & Any repeat transfer & No enabled transition & No second transfer is permitted by this model.\\
\bottomrule
\end{longtable}

Start with $p=100$, $b=0$, $r=0$, and an agreed deposit $d=40$. A valid deposit
produces $(p,b,r)=(60,40,0)$ and state \texttt{HELD}. If delivery is accepted before
expiry, release produces $(60,0,40)$ and state \texttt{PAID}. If expiry occurs
first, the refund produces $(100,0,0)$ and state \texttt{REFUNDED}. The same deposit
has two possible future paths, determined by events and their timing.

A reader can now check a conservation property without trusting a compiler. Define
$W=p+b+r$, the total modeled resource. For the deposit transition,
\begin{equation}
 W'=(p-d)+d+r=p+r=W,
 \label{eq:escrow-deposit}
\end{equation}
because the initial contract holding is zero. For release,
$W'=p+0+(r+b)=W$. For refund, $W'=(p+b)+0+r=W$. Each permitted transition therefore
preserves the total. Since the initial holdings are nonnegative and the deposit
guard requires sufficient funds, these transitions also preserve nonnegativity.
This is a local derivation about the stated miniature model.

The terminal states prevent a second transfer. After release, no transition from
\texttt{PAID} can release the same balance again. A Java implementation that
updates balances but forgets the state change could violate that property under
a repeated call. This is why the state transition is part of the contract, not
just a user-interface label. An idempotency mechanism can support repeated
requests, but must be specified consistently with the state model.

Time is equally substantive. The table uses $t<\tau$ for release and
$t\geq\tau$ for expiry. At exactly $\tau$, only expiry is enabled. If the intended
agreement allowed delivery at the deadline, the guards would need a different
boundary or an explicit ordering rule. A test performed one minute before expiry
cannot establish the exact-deadline behavior. The Stipula literature's event
ordering is therefore a modeling choice to understand, not a universal answer
for every contract \citep{stipula2021}.

The model also exposes a dependency hidden by ordinary prose: what establishes
accepted delivery? It might be an authorized person's recorded acknowledgment,
an external system event or a disputed factual assessment. The transition table
assumes that this predicate is supplied correctly. Conservation and no-double-
transfer proofs do not establish that assumption. If acceptance can be forged,
the formal transfer properties may hold while the wrong person receives the
resource.

The distinction between safety and progress is now visible. Conservation is a
safety property: no permitted step creates or destroys the modeled resource.
The ability to reach \texttt{PAID} or \texttt{REFUNDED} is a progress question.
If the platform never processes the expiry event and delivery is never accepted,
the resource can remain held. Proving eventual refund requires assumptions about
time advancement and event execution, not merely the transition table's arithmetic.
The contract-liquidity literature concerns such availability of actions and
resources; it should not be confused with market liquidity \citep{liquidity2023}.

This miniature also clarifies what a Java/JML translation would need to preserve.
A precondition specifies the caller, state and guard. A postcondition specifies
the new balances and state. An invariant states nonnegative holdings and any
conserved total. The verifier reasons about all allowed calls under those
conditions. It must also account for the event mechanism and any loop abstraction
used by the translation. The inspected Stipula-to-Java work provides a relevant
route, with the restrictions already described \citep{stipulakey2025}.

For LegalMath, the corresponding benefit is a disciplined way to model review,
consent and obligation lifecycles. A report awaiting review is not an approved
report; a withdrawn consent is not an active generation; an opened duty is not
fulfilled merely because a document was eventually delivered late. Explicit
states and transitions make these distinctions executable and testable.

The correspondence is methodological rather than a claim that every circular is
an escrow contract. Regulatory obligations can apply even when a party never
agreed to a private contract, and actual events must be recorded even if a formal
platform would reject them. A bank's monitoring system therefore separates the
normative question of what should occur from the factual record of what did
occur. The state model helps express that separation only when it is designed
for the intended purpose.

A practical specification exercise follows directly. Before implementing a new
history-dependent control, write its states, actors, events, guards and effects;
execute one normal path, one boundary path and one violating observation by hand;
identify the facts that the model assumes; and state which invariants and progress
properties matter. Then choose a language or verifier that supports those
properties. This sequence gives Stipula, Symboleo, eFLINT or a restricted custom
monitor a concrete role rather than selecting a language because its name occurs
frequently in the literature.

\section{What contract verification does and does not buy}
\label{sec:languages-contractverification}
Once a contract has a transition semantics, tools can ask whether certain
states are reachable, whether a resource can become trapped, or whether two
contracts have corresponding observable behavior. Stipula's legal-bisimulation
work gives a formal account of selected observable equivalences
\citep[behavioral equivalence section]{stipula2021}. A bisimulation relates
states so that an observable step on one side can be matched by a corresponding
step on the other. It is a statement about modeled behavior under that
definition of observation.

For the bank, an analogous question might ask whether a revised consent
implementation preserves the reviewed behavior for grants, withdrawals and
transactions. The observation set must include what matters: subject, category,
time, permitted use and recorded violations. If the comparison ignores
withdrawal, equivalence of the remaining observations would not establish the
property the bank needs. Formal observation boundaries deserve the same
scrutiny as ordinary input-domain restrictions.

\citet{stipulakey2025} demonstrate a translation from Stipula contracts into
Java with JML specifications and verification using KeY. JML expresses properties
of Java behavior, such as preconditions, postconditions and invariants; KeY
reasons about whether the Java program satisfies them. This is especially
relevant to our Java delivery target because it illustrates a route from a
specialized contract language into a general-purpose verification environment.

The inspected work has explicit restrictions. Its event conditions are modeled
symbolically, and the automated encoding keeps certain loop inputs constant
across iterations. Its examples establish feasibility for the examined models,
not complete fidelity to arbitrary calendars or external event streams. The
official translator source was inspected in the earlier project review; no KeY
proof was replayed. Adopting a similar route would require a precise mapping
from our RuleIR and event profile to the Java/JML semantics, including the
unknown and conflict results that ordinary Boolean contracts omit.

General verification cannot be promised for every expressive contract language.
\citet{reachability2026} establish undecidability results for clause reachability
in Stipula fragments and identify a decidable fragment under particular
restrictions. The restrictions concern disjoint function/event source states
and immediate events; they should not be confused with the different
disjoint-cycle condition used in the Java/JML translation work. A shared word
in two papers does not identify a shared theorem.

The practical consequence is to define bounded questions and supported
fragments. A tool can report that no bad state was found up to a specified
depth, or prove a property under a restricted transition model. It cannot
turn a finite search into a complete reachability decision for a language
whose general problem is undecidable. The report must retain the bounds so a
reviewer can understand which behaviors were left unexplored.

Stipula's amendment literature is also relevant. A contract can change its
declarations, state and permissible behavior through a specified amendment
mechanism \citep[secs.~3--5]{amending2023}. For regulatory controls, an amendment
must similarly identify whether existing obligations keep their opening source
version or migrate to a new one. Hot-replacing code without a state-migration
interpretation can change the treatment of historical events. The versioned
control and its transition obligations need separate approval.

\section{Other normative languages clarify different distinctions}
\label{sec:languages-other}
Symboleo models legal contracts using concepts such as obligations and powers,
with lifecycle distinctions that are easy to lose in a Boolean implementation
\citep[secs.~II--IV]{symboleo2020}. Creating a duty is not always the same as
activating it. A power can change a normative relationship, for example by
suspending or terminating an obligation. These distinctions are useful when a
bank workflow has review, activation, withdrawal and corrective actions with
different legal effects.

The retained Symboleo paper and software inspection do not provide a complete
implemented axiom set for LegalMath. Some supporting references in the early
paper are not fully available in the retained edition. We can use the lifecycle
questions to improve requirements, while treating any direct engine adoption
as a separate interface and semantics audit. A familiar conceptual diagram is
not a substitute for an executable contract.

eFLINT separates facts, actions and duties in a normative specification
\citep[secs.~3--5]{eflint2020}. A particularly useful distinction is that an
impermissible action may still occur, change state and create a violation.
This matches the difference between a bank's pre-transaction advice and its
post-event audit. A system that refuses to represent an unauthorized observed
action cannot explain what happened after a control failure.

The eFLINT literature also shows why language versions matter. The later
stable-model semantics provides a particular formal grounding and can produce
no model or several models \citep[translation and evaluation sections]{eflintstable2025}.
Those outcomes carry information; selecting one model silently would conceal
ambiguity or inconsistency. The newer language survey discusses implementations
whose semantic guarantees differ \citep{eflint2026}. A proof or test result for
one version should not be attributed to every interpreter bearing the same name.

s(LAW) combines legal formalization with logic programming that can explain
positive and negative outcomes and explore possible models
\citep[secs.~2--5]{slaw2024}. This makes assumptions visible in a way that a
single opaque classification may not. However, a model containing a hypothetical
fact does not establish that the fact is true for the client. Counterfactual
reasoning should propose what would change an answer, while the evidence
service establishes whether the corresponding condition actually holds.

These methods suggest an interface in which normative classification,
state transition and factual evidence are separate inputs to a decision. The
implementation should reject an unsupported duty kind explicitly. A maintenance
obligation, which must hold throughout an interval, cannot be implemented by
checking that it was true once. A non-preemptive achievement obligation cannot
be discharged by a performance before it came into force. The taxonomy in
\citet[sec.~3]{normative2016} helps identify these choices; the precise profile
still needs a source-grounded interpretation.

\section{Decision tables and business-rule engines}
\label{sec:languages-tables}
Decision tables remain useful when their rows and hit policy have a precise
meaning. A table can show, for example, that a known discount exception changes
the selected trigger while an unknown exception leaves a material question.
The reader must know what happens when several rows match: choose the first,
require a unique row, collect results or apply a priority. DMN makes such hit
policies explicit \citep[sec.~8.2.11]{dmn2024}.

That explicitness is the valuable part. A spreadsheet whose row order silently
resolves competing legal exceptions is harder to review because its legal
priority is hidden in presentation. Changing the sort order can change the
answer. A unique-hit table can reject overlap, but it cannot determine whether
the missing row represents an omitted legal situation. Independent source
coverage and boundary cases remain necessary.

Drools, OPA and other rule engines can execute policy logic and integrate with
enterprise applications. The question for this product is not whether they
are capable software. It is whether the selected engine's semantics, evidence
model, explanation behavior and deployment interface match the adopted
specification. Unknowns, conflicts, exception priorities and event ordering
need explicit mappings. A backend name is not an assurance claim.

For the prototype, a small RuleIR is easier to specify and compare than a
large engine with many unsupported features. Other engines can serve as
independent reference encodings for a supported fragment after that fragment
is defined. The comparison must include the behavior that ordinary Boolean
examples leave out. Engine diversity without semantic alignment can produce
many discrepancies that reflect adapter choices rather than flaws in the
legal interpretation.

\section{A solver can improve reasoning after a translation}
\label{sec:languages-neurosymbolic}
A neural-symbolic system separates tasks that language models and formal
engines perform differently. A language model proposes a formal representation
of prose. A symbolic engine evaluates or reasons about that representation.
The engine can give a reproducible answer to a precise logical question and
can expose inconsistency, unsupported syntax or a counterexample. The
translation that supplied its premises remains a separate source of error.

Logic-LM uses language-model formulation followed by an appropriate solver and
feedback-driven refinement \citep[sec.~3 and prompt appendices]{logiclm2023}.
LINC combines language-model formalization with first-order theorem proving
and analyzes failures at the translation and reasoning boundary
\citep[method and error analysis]{linc2023}. These methods support the design
choice that a model should not be asked to simulate every logical calculation
in prose. They do not support treating successful solver execution as evidence
that a legal requirement was fully captured.

A syntax repair illustrates the benefit and the limit. Suppose a candidate
uses an undeclared symbol in the gift rule. The parser identifies the symbol
and rejects the encoding. A repair prompt can ask for the missing declaration
or corrected reference while preserving the intended interpretation. The
repaired candidate is then validated again. This process may solve a concrete
representation error. It does not justify adding a new legal assumption
merely because doing so makes the formula easier to satisfy.

ARc develops policy-model creation, human vetting, structured inspection and
symbolic checking in a combined system \citep[sec.~3 and appendix A]{arc2025}.
Its policy creation procedure divides a document into units, proposes types,
variables and constraints, and composes them into a model. Its vetting
mechanisms include linting, template-based English inspection and testing.
These are relevant to our product because they make the proposed model
available for correction before individual answers rely on it.

Composition creates its own interpretation problem. Two units may use the same
word for different legal concepts, or different words for the same concept.
ARc uses embedding-based clustering in its variable-unification procedure.
For our high-stakes specification, such a merge should be a proposal with
provenance and validation. Merging ``offer discount'' with ``component
discount'' because their descriptions are similar could remove the distinction
that g04 was designed to test. A similarity score cannot establish identity
of legal concepts.

ARc's answer-verification procedure also uses redundant translations and
symbolic agreement checks. Its reported finding categories distinguish
ambiguous translation, impossible premises, invalid conclusions, entailed
conclusions and satisfiable but unentailed claims. The distinction is useful:
if a model permits both an outcome and its opposite under incomplete premises,
the system can show scenarios for each rather than invent a definitive answer.
For LegalMath, these scenarios become review questions tied to the candidate
set and the source context.

The word ``soundness'' in ARc's empirical evaluation needs careful translation
into a product claim. The paper defines an empirical quantity based on false
positives across requests; it is not a theorem that the English policy was
formalized correctly. The denominator differs from precision among approvals.
The manuscript therefore borrows the inspection and disagreement mechanisms
without importing a claimed legal-error probability. Chapter~\ref{ch:evaluation}
defines the denominators our own evaluation would need.

The tax-reasoning work of \citet{jurayj2026} further illustrates the importance
of task boundaries. Supplying gold rules, retrieving annotated cases and
extracting facts are different assistance conditions from generating rules
from a new circular. A comparison must account for those inputs and the human
effort that produced them. A system given a reviewed rule library should not
be presented as evidence that the same system can independently interpret
unreviewed prose.

\section{Roundtrip verification and targeted repair}
\label{sec:languages-roundtrip}
Roundtrip verification asks whether a formal representation survives a journey
through another representation. In the method of \citet{roundtrip2026}, source
English is formalized, the formula is translated back into English, and the
reconstructed English is formalized again. A solver compares the first and
last formulas. If they differ, diagnosis identifies a likely translation
stage to repair, and dependent stages are regenerated.

This is attractive for our user interface because the intermediate English
can make a formal candidate easier to inspect. Suppose the original formula
contains both product-link routes but the back-translation mentions only a
specific product. Re-formalizing that sentence may lose the product-type
route. The resulting counterexample exposes a defect in the explanation or
translation chain. A reviewer should not be shown that explanation as if it
faithfully described the formal rule.

The same construction can miss a shared omission. If the first formalization
already omits the product-type route, a faithful back-translation of that
incorrect formula and a faithful re-formalization will agree. The roundtrip
then passes. What survived is the first formula, including its error relative
to the source. The paper explicitly recognizes this limit; our ensemble must
retain a source-grounding path that does not simply inspect the same
back-translation.

Repair therefore needs a named target. If the formal rule is justified but
the explanation drops a condition, repair the explanation and its dependent
roundtrip outputs. If the original rule omits a source condition, create a
new interpretation candidate and revalidate its downstream artifacts. If the
two formulas use incompatible vocabularies, repair the mapping before claiming
the formulas disagree about one question. A generic ``make these consistent''
instruction does not distinguish those cases.

The stopping rule is equally important. A roundtrip procedure can run for a
bounded number of attempts and still end in disagreement. The product should
report the original and revised formulas, the counterexample, diagnosis,
actions attempted and the remaining issue. It should not keep repairing
until some pair agrees and then discard the path that made the original
source objection visible. The ensemble controller extends this principle
from one translation chain to several methods and issue types.

An explanation generated by deterministic templates provides a useful
additional path. It is constrained by the formal syntax and can be checked
for completeness of its displayed operands. It may be less elegant than
free prose, but it offers a stable view of what the program actually says.
A freely generated explanation can still help readers if it is compared
against that view and does not replace it as the sole account of the rule.

\section{Evidence spans and distant qualifications}
\label{sec:languages-evidence}
An answer supported by a source span is easier to review than an unsupported
assertion, but the span itself may be insufficient. A qualification can appear
in a footnote, definition or distant section. ContractNLI studies document
inference with evidence annotation and highlights challenges involving
references and exceptions \citep[task, annotation and error-analysis sections]{contractnli2021}.
For our product, the relevant lesson is to retain the supporting and limiting
passages needed for the inference, rather than a single convenient quotation.

The distinction between ``not mentioned'' and ``false'' is particularly useful.
If a retained source does not settle the definition of a fee refund, the
absence should remain visible. A retrieval engine returning no matching
paragraph does not establish that the exception excludes refunds. It may have
searched the wrong document, missed a synonym or lacked the referenced source.
The evidence report should distinguish unsuccessful retrieval from reviewed
absence within a specified source set.

Source spans also require stable identity. Character offsets depend on the
extraction version, while page numbers may not survive a changed layout.
Retain raw-byte identity, derivative identity, exact span text or hash, and
human-readable location. A reference to paragraph 10 without a source version
can become misleading after an amendment. A matching text hash without its
source context may identify a sentence copied into a secondary summary.

The formalization boundary should distinguish source claims from supplied
facts. A sentence defining a gift exception is part of the normative context.
A bank record establishing that a fee was paid is factual evidence. A
reviewer's classification relating the record to the exception is an
interpretive application. Combining all three into an undifferentiated
``evidence'' list makes it harder to see where disagreement actually arises.

Recent faithfulness work analyzes how language models and formal solvers can
diverge when implicit assumptions enter formalization \citep{limits2026}.
The retained paper also has annotation and denominator questions that prevent
us from transferring a general accuracy figure. Its examples nevertheless
reinforce a practical requirement: an implicit assumption that changes the
answer should become an explicit, challengeable premise. The system must
show where a model supplied information that the source packet did not.

\section{Testing paths, dates and inconsistent constraints}
\label{sec:languages-analysis}
Other methods contribute different checks once a formal model exists.
CUTECat uses concolic execution to explore paths in Catala programs
\citep[secs.~2--5]{cutecat2025}. Concolic execution combines a concrete run
with symbolic constraints describing the path taken. By changing a path
condition, the tool seeks an input that exercises another branch. For our
gift example, this is analogous to moving from a named-product case to the
case where only the product-type route can determine the answer.

Path-directed generation can expose an exception branch that ordinary
examples never visit. It cannot visit a branch that the formal program
omitted entirely. Independent source-derived cases remain necessary. The
reported costs on larger legal examples also show why a prototype needs
explicit resource limits; small-program speed should not be treated as a
promise that every circular can be exhaustively explored.

Date arithmetic is another place where a precise program can embody a wrong
choice. Adding a month to a date near the end of a month may require a
rounding convention. The choice can affect deadlines. \citet{dates2024}
formalize date arithmetic and analyze ambiguity in legal applications. The
transfer to this project is to make calendar rules, boundary conventions and
their source basis explicit. A verified date library does not decide which
rounding convention a particular provision requires.

ContractCheck analyzes consistency and executability of structured contract
models \citep[secs.~3--8]{contractcheck2025}. A satisfying assignment shows
that a modeled set of conditions can jointly hold. An unsatisfiable set can
help locate conflicting clauses or assumptions. Neither result proves that
the original legal document has been translated faithfully. The study's
manual translation is an important part of its task boundary.

An unsatisfiable core is a useful diagnostic: a subset of constraints that
already suffices for inconsistency. It is not necessarily the smallest such
subset unless a minimization procedure establishes that property. For a
reviewer, a small and source-linked explanation is valuable, but the product
should label a returned core accurately. It must also distinguish an
encoding inconsistency from legitimate defeasible conflict that the chosen
logic was intended to resolve.

These analyses fit into the ensemble as specialized questions. Does this
representation have a hidden path? Does a date convention alter a boundary
case? Are the supplied assumptions jointly possible? Does a revised rule
preserve the reviewed behavior? Their evidence can challenge a candidate
without claiming to replace the source interpretation process.

\section{A synthesis for the first implementation}
\label{sec:languages-synthesis}
The immediate product should keep a small executable core whose semantics are
fully specified. Around it, the interpretation record should be richer than
the executable language: it needs alternative readings, source versions,
definitions, unresolved qualitative judgments and supporting arguments.
Controlled templates give reviewers a readable view of formal structure.
Specialized analyses can then inspect declared fragments through adapters.

This choice avoids forcing every legal question into the first compiler.
A paragraph requiring judgment can become a reviewed assessment input or
manual procedure, with explicit evidence and responsibility. A temporal duty
can use an event profile whose limits are stated. A disputed definition can
remain in the interpretation set until review resolves it. The system can
still generate Java for a reviewed subset while keeping the broader
assessment visibly incomplete.

Every adapter needs a semantic contract. It declares supported constructs,
result categories, unknown and conflict behavior, vocabulary mapping,
resource bounds and how a returned witness is replayed. A paper's name or
a repository's existence does not supply that contract. A failed or
unsupported analysis is reported as such and cannot be counted as an
additional agreeing ensemble member.

The literature thus supplies complementary mechanisms: source-linked
authoring, precise exceptions, contextual alternatives, normative state,
symbolic consequence, path-directed testing and translation repair. The
new contribution proposed here is the controller that combines these
mechanisms without allowing one successful check to conceal another
unresolved question. Before specifying that controller, we need to explain
how competing interpretations are generated and defended. That is the task
of the next chapter.
